\section{Experimental setup}%
\label{sec:experiments}

\noindent\textbf{Benchmarks} We use a thoroughly validated, reproducible subset of \textbf{200 instances} from the \kbenchsyz{} benchmark~\citep{kgym} of $279$ Linux kernel crashes found by the Syzkaller fuzzer~\citep{syzkaller}.
Each instance in the benchmark consists of (1) a reproducer file, containing the user-space program that triggers the crash, (2) the ground-truth commit that fixed the bug, 
and (3) the crash report at the parent commit of the fix commit (we run all tools at this parent commit).
To show generalizability, we also evaluated \cresearcher{} on $10$ recent crashes of an open-source multimedia software, FFmpeg~\citep{ffmpeg}.
More details about the \kbenchsyz{} benchmark and the \ffmpeg{} dataset are in Appendix~\ref{app:details} and Appendix~\ref{app:ffmpeg} respectively.

\noindent\textbf{Evaluation metrics}
We compute Pass@k~(\textbf{P@$k$}) defined as $\textbf{P@$k$} = 1$, if applying at least one of the $k$ candidate patches generated by the tool prevents the crash, 
or $\textbf{P@$k$} = 0$ otherwise. We report \textbf{(1)} \textbf{C}rash \textbf{R}esolution \textbf{R}ate~(\textbf{CRR}) which is average \textbf{P}@$k$, \textbf{(2)} average \textbf{recall}, i.e., the fraction of files modified in the ground-truth commit (\emph{the ground-truth buggy files}) in the set of files edited by the agent, averaged over the $k$ candidate patches, and 
\textbf{(3)} the percentage of candidate patches where \textbf{All}, \textbf{Any} or \textbf{None} of the ground-truth buggy files are edited. 
When a tool does not produce a patch (e.g., it runs out of LLM call budget), the set of edited files is assumed to be empty. All the metrics are averaged over the $200$ instances in the benchmark. 

\noindent\textbf{Baselines}
We evaluate \textbf{\cresearcher{} in the \emph{unassisted} setting} (i.e., the ground-truth buggy files that are part of the fix commits are not divulged to the tool) and compare it against the following baselines:
\textbf{(1)} \textbf{o1}~\citep{o1} \textbf{and GPT-4o}~\citep{gpt4o} \textbf{in the \emph{assisted} setting}, i.e., we give the ground-truth files that are part of the fix commits and the crash report as input. We prompt the model to generate a hypothesis about the crash's cause and a patch.
\textbf{(2)} \textbf{o1 and GPT-4o in the \emph{stack context} setting}, where we give the contents of the files mentioned in the crash report as input besides the crash report. 
\textbf{(3)} \textbf{SWE-agent 1.0}~\citep{yang2024sweagentagentcomputerinterfacesenable}, a SOTA coding agent on the SWE-bench benchmark, \textbf{in the \emph{unassisted} setting}.
For fairness, we add a Linux kernel-specific example trajectory and background about the kernel to its prompts.
We sample $k$ (for Pass@$k$) SWE-agent trajectories independently using a temperature of $0.6$.
\textbf{(4)} \textbf{Agentless}~\citep{xia2024agentlessdemystifyingllmbasedsoftware}, another top coding agent on the SWE-bench benchmark that uses a fixed workflow of localization followed by repair, \textbf{in the \emph{unassisted} setting}. We add background about the kernel to its prompts and sample $k$ (for Pass@$k$) patches using the same LLM (GPT-4o) as \cresearcher{} and SWE-agent. 
\textbf{(5)} \textbf{CrashFixer}~\citep{crashfixer}, state-of-the-art agent for Linux kernel crash resolution, \textbf{in the \textit{assisted} setting}, as it directly uses the ground-truth files to generate patches. If a patch fails to build or crashes with the ground-truth reproducer, it iteratively refines it using the respective error messages. We report their results with the Gemini 1.5-Pro-002 model (v. 2024-09-24) and more recent results with the Gemini 2.5-Pro model (v. 2025-06-17).

\noindent\textbf{Implementation, hyperparameters}
We employ \textbf{GPT-4o} (v. 2024-08-06) for \cresearcher{} and for the competing tools. We also experiment with \textbf{o1} (v. 2024-12-17) in the \synthesisphase{} phase of \cresearcher{}, and,
to show that \cresearcher{}'s design is not tied to the OpenAI family of models, the newer \textbf{Gemini 2.5-Flash} (v. 2025-06-17) for both the \analysisphase{} and \synthesisphase{} phases.
All our experiments have a context length limit of $50$K tokens.
In the \analysisphase{} phase, we use a \textbf{temperature} of $0.6$ and independently sample $k$ trajectories.
For the \synthesisphase{} phase, we sample with increasing temperatures ($0$, $0.3$, $0.6$) until the agent produces a correctly-formatted patch, with a maximum of $3$ attempts.
We allow all tools a budget of at most \textbf{\maxcalls{}} LLM calls to generate a single patch.
Please refer to Appendix~\ref{app:details} for details about the \textbf{crash reproduction setup}, the \textbf{prompts}, the \textbf{compute resources} and the \textbf{configurations} for \cresearcher{} and the baselines.
